\subsection{Executed cover: applying the procedure}\label{sec:cover-study}
Two fixed-weight models, Qwen3:8B and Llama3.1:8B, each run 300 synthetic ticket fixtures per condition: 120 training, 60 validation and 120 test, balanced across six actions. The independent test unit is the fixture: 120 per model, paired across six conditions. Request phrasings recur across splits.

\paragraph{Execution and observations}
LangGraph~\cite{langchain2025langgraph} lets each model choose checks, actions, objects, recovery and finish. Seven local MCP Streamable HTTP destinations use TLS~1.3/HTTP/2, a common \path{/mcp} path, and 4,096/16,384-byte request/response padding quanta. A frozen schedule places one or two decoys among early calls. Policies are no cover, proxy discard (computes on a disposable copy), read cover, hidden isolated execution (on a retained clone), live writes (an adverse state-pressure control on the primary fixture), and visible isolation (marked cover responses reach the planner). The observer receives destinations and bidirectional ciphertext counts in supplied episode/call windows. Goal completion checks the requested action/object, prerequisites and target postcondition; whole-state integrity also checks other orders, shared quota and the ledger. Intervals resample paired fixtures by action, conditional on fitted attacks.

\begin{table*}[t]\centering\small
\caption{Executed cover on the same held-out executions (120 paired fixtures per model and policy; chance 20/120). \emph{Size view}: Byte is the validation-selected size/destination learner, Last the last-destination rule, and Best the validation-selected size-only attack among learners and position rules. \emph{Unmatched} (client-receipt call times): Strip removes calls flagged by the label-free gap rule and applies the frozen no-cover classifier; Jit.\ is the refitted strip attack under 200\,ms half-normal jitter (mean of five draws). \emph{Timing-matched} (gap and duration drawn from no-cover history): Strip uses the refitted gap detector; Adapt.\ is the validation-selected attack after adding the context detector and joint decoding; +Lat is added latency as a percentage of no-cover wall time. \emph{Execution}: Goal, whole-state integrity (State) and completed tasks with state errors. Call bytes are 1.28--1.35$\times$ under cover. Supplements S9--S11 give all attacks, intervals and tests.}
\label{tab:cover-main}
\setlength{\tabcolsep}{3.2pt}
\begin{tabularx}{\textwidth}{@{}lYrrrrrrrrrrr@{}}\toprule
 & & \multicolumn{3}{c}{Size view} & \multicolumn{2}{c}{Unmatched} & \multicolumn{3}{c}{Timing-matched} & \multicolumn{3}{c}{Execution}\\\cmidrule(lr){3-5}\cmidrule(lr){6-7}\cmidrule(lr){8-10}\cmidrule(l){11-13}
Model & Cover policy & Byte & Last & Best & Strip & Jit. & Strip & Adapt. & +Lat\% & Goal & State & Goal+err \\\midrule
Qwen3:8B & No cover & 118 & 116 & 118 & -- & -- & -- & -- & -- & 114 & 116 & 0 \\
Qwen3:8B & Proxy discard & 62 & 77 & 77 & 115 & 104 & 43 & 76 & 12.4 & 114 & 116 & 0 \\
Qwen3:8B & Read cover & 62 & 77 & 77 & 114 & 103 & 50 & 84 & 12.4 & 114 & 115 & 0 \\
Qwen3:8B & Isolated, hidden & 61 & 77 & 77 & 114 & 106 & 57 & 84 & 12.8 & 114 & 115 & 0 \\
Qwen3:8B & Live write & 62 & 75 & 75 & 112 & 102 & 54 & 82 & 12.7 & 102 & 7 & 95 \\
Qwen3:8B & Isolated, visible & 63 & 77 & 77 & 111 & 98 & 49 & 86 & 12.3 & 113 & 114 & 1 \\
\midrule
Llama3.1:8B & No cover & 113 & 79 & 113 & -- & -- & -- & -- & -- & 82 & 69 & 13 \\
Llama3.1:8B & Proxy discard & 79 & 67 & 79 & 116 & 115 & 59 & 95 & 16.4 & 80 & 68 & 13 \\
Llama3.1:8B & Read cover & 72 & 67 & 72 & 112 & 110 & 59 & 89 & 15.8 & 84 & 71 & 14 \\
Llama3.1:8B & Isolated, hidden & 82 & 67 & 82 & 112 & 110 & 55 & 97 & 15.7 & 83 & 70 & 13 \\
Llama3.1:8B & Live write & 80 & 65 & 65 & 112 & 110 & 66 & 92 & 16.4 & 76 & 3 & 73 \\
Llama3.1:8B & Isolated, visible & 81 & 68 & 81 & 107 & 105 & 62 & 87 & 16.2 & 83 & 73 & 11 \\
\bottomrule\end{tabularx}
\end{table*}


\paragraph{Step 1: size-only evaluation, with simple rules}
The validation-selected learner scores 61--63/120 for Qwen and 72--82 for Llama under cover, against 118 and 113 without it. A learner alone overstates the protection. The last-destination rule reaches 75--77 for Qwen, and the validation-selected best of learners and position rules gives 75--77 and 65--82 (Best column). This is the cover's size-only protection level, at 1.28--1.35$\times$ call bytes.

\paragraph{Step 2: a removal attack built on the execution mechanism}
Honest tool calls follow a model decision; the controller dispatches cover directly after a tool response. We add \emph{client-receipt timing}, the start and end of each call at the Agent's network endpoint. An unsupervised two-component mixture on unlabeled training gaps sets a threshold of 0.19\,s (Qwen) and 0.14\,s (Llama). On test it flags all 1,820 decoys and none of 6,741 honest calls. Removing them and applying the frozen no-cover classifier restores 107--116/120; all ten gains over the byte learner survive Holm correction ($p\leq6.2\times10^{-6}$). The margin is large: on validation, cover gaps never exceed 56\,ms and honest gaps are never below 418\,ms. With independent half-normal jitter on every timestamp, a refitted detector keeps 98--115/120 at $\sigma=200$\,ms. A generic LR/RF given the same timing features does not improve on the size view (no primary test rejects). Timing helps the attacker only through a removal step that models the execution, as in WF where attack design rather than feature availability decided defenses~\cite{dyer2012peekaboo,rahman2020tiktok,mathews2023critical}.

\paragraph{Step 3: timing-matched dispatch, and attacks that adapt to it}
Adaptive padding and WTF-PAD draw dummy timing from the protected traffic's own distribution~\cite{shmatikov2006adaptive,juarez2016wtfpad}. The Agent analogue delays each decoy by a gap, and extends it to a duration, drawn from the defender's no-cover history. We simulate it on the collected receipts. Hidden cover never reaches the planner, so its decisions are unchanged. Against the refitted gap detector, matching works: precision falls to .46--.57 and stripping drops to 43--66/120. The earlier pre-registered criterion, strip within ten points of the size view, is nonetheless not met in seven of ten cells, because stripping falls further below it.

A second frozen analysis tests attackers that adapt to the matched cover. A context detector compares each call with its episode, position and destination, and joint decoding scores every removal of one or two early calls together with the classifier's confidence. The context detector separates matched decoys at AUROC .93--.94. The validation-selected adaptive attack exceeds the size-only best in nine of ten cells (Adapt.\ column): by $-1$ to 9 tasks for Qwen and 6 to 27 for Llama. Two Llama cells survive Holm correction (read, +17, $p=.030$; live write, +27, $p=.002$). In live write, validation picks the last-destination rule (65) although the byte learner scores 80; against it the gain is 12. The pre-registered rule for declaring the countermeasure defeated, at least 12 tasks with Holm $p<.05$ in five of ten cells, is not met. Drawing gaps by call position instead leaves the same two significant Llama cells.

\paragraph{Why matched timing is not enough}
A post-hoc ablation of the context detector shows two separate signals. Without any timing input, destination and position context alone separate matched decoys at AUROC .83--.85; timing inputs alone reach .80--.82. A decoy schedule drawn independently of the planner's destination sequence is visible even when its timing is matched. The timing signal remains because the detector compares a gap with the rest of its own episode, which neither a pooled nor a position-wise draw reproduces. Cover generated outside the model loop must therefore imitate the planner's joint sequence and timing behavior. Matched marginals cost 12--17\% added latency and are not sufficient. These are receipt-level counterfactuals on loopback (Supplement S11).

\paragraph{Business outcomes remain a separate constraint}
Proxy, read and hidden-isolated cover keep Qwen state integrity at 115--116/120. Adverse live writes reduce it to 7/120, and Llama falls from 69 to 3. Of 102 Qwen and 76 Llama completed live-write tasks, 95 and 73 have collateral state errors, so completion-only reporting would pass them. This is a deliberately adverse control, not an estimate of production harm. Visible cover responses raise mean prompt tokens by 18.1--18.9\%, with no significant endpoint change (Supplement S8). Schedule and timing determine what the observer learns; semantics determine what the application pays.
